\chapter[Bibelverfälschung]{Bibelverfälschung}
\quellseite{7}
Die neuere Theologie in Deutschland wie in America ist ja schon längst so gut wie allgemein vom alten Bibelglauben abgefallen und hält die Schrift nicht mehr für Gottes Wort. Diesen ihren Unglauben will sie nun auch immer mehr unter das Volk bringen und versucht es auf mancherlei Weise. Das Neueste ist, daß sie nun auch eine neue Bibel an die Stelle der alten setzen möchte. In einer verbreiteten deutschen kirchlichen Zeitschrift, die sich die „Christliche Welt“ nennt, werden „Worte Jesu, wie sie heute klingen könnten“, gedruckt. Der Verfasser, Walter Classen, will dem Volke unserer Tage den Herrn Christum zeigen, wie die moderne Theologie ihn sieht; und weil Christi Worte nach der Meinung dieser Theologie so unmodern geworden sind, sollen sie nun recht modern gemacht werden. Und dabei werden dann mit freveler Hand die Worte des Heilandes ganz umgestaltet nach Wortlaut und Inhalt. Das Wort: „Selig

\quellseite{8}
sind die Barmherzigen, denn sie werden Barmherzigkeit erlangen“, Matth. 5, 7., lautet nun so: „Ich preise die Hülfreich-Thätigen, denn jeder hilft ihnen gern.“ Die andere Seligpreisung: „Selig sind, die reines Herzens sind, denn sie werden Gott schauen“, Matth. 5, 8., wird so wiedergegeben: „Ich preise, die eine reine Gesinnung haben, denn sie können Gott sehen in der Welt.“ Und so geht es fort, und mit frecher Hand wird das Wort Christi verdreht. — Gerade so macht es in America die in Chicago erscheinende Zeitschrift The Biblical World, die von den theologischen Professoren der dortigen großen Universität herausgegeben wird. In jeder Nummer des letzten Jahrgangs nahm sie einen Bibelspruch her, stellte die verschiedenen Uebersetzungen desselben zusammen und brachte dann zum Schluß eine eigene Fassung. Das Ganze nennt sie “A Study in modernizing the English Bible”. Wie sie dabei verfährt und heilige Gottesworte in den Koth ihrer falschen menschlichen Meinung tritt, mögen zwei Beispiele zeigen. Der Spruch Ps. 19, 13., den die altenglische Bibel so übersetzt: “Who can understand his errors? Cleanse Thou me from secret faults”, wird so wiedergegeben: “No one can avoid errors due to human weakness and ignorance; do Thou therefore not hold me responsible for offenses unconsciously committed.” Die Ermahnung St. Pauli Röm. 12, 1. 2. lautet in der englischen Bibel so: “I beseech you, therefore, brethren, by the mercies of God, that ye present your bodies a living sacrifice, holy, acceptable unto God, which is your reasonable service. And be not conformed to this world, but be ye transformed by the renewing of your mind, that ye may prove what is that good, and acceptable, and perfect will of God.” Die Biblical World aber modelt dieses Wort so um: “Fellow Christians, God is very loving and good to us; and His will is the only true guide to life. Therefore free yourselves from sinful practices, be pure and noble, think and act according to the gospel teaching; this-is the kind of worship which God wishes from you.” L. F.
